\section{Arithmetic inputs and exact normalizations}\label{sec:classical}

The two comparisons will be carried out at a fixed odd prime. This
section explains the reduction to that problem, constructs the
quadratic fields used in the pair comparison, and records the real
arithmetic identity that converts a Heegner-point index into a sum of
BSD discrepancies. None of these reductions asserts an odd-primary
leading-term equality.

\subsection{The two low-corank inputs}

For a number field $M$, the compact local Kummer condition for the
Tate lattice $T_pE$ is the image of the completed point module
\[
 E(M_v)^{\wedge}_p=\varprojlim_n E(M_v)/p^nE(M_v)
       \longrightarrow H^1(M_v,T_pE).
\]
For discrete coefficients, the local condition is the image of
$E(M_v)\otimes\Q_p/\Z_p$ in $H^1(M_v,E[p^\infty])$. The full Selmer
group uses these conditions at every place. The Kummer exact sequence
is
\begin{equation}\label{eq:kummer-exact}
 0\longrightarrow E(\Q)\otimes\Q_p/\Z_p
 \longrightarrow\Sel_{p^\infty}(E/\Q)
 \longrightarrow\Sha(E/\Q)[p^\infty]\longrightarrow0.
\end{equation}

\begin{theorem}[Unrestricted low-corank converse]\label{classical:low-corank}
For every elliptic curve $A/\Q$, every prime $q$, and $u\in\{0,1\}$,
the equality $s_q(A)=u$ implies
\[
 a(A)=r(A)=u,\qquad \#\Sha(A/\Q)<\infty.
\]
\end{theorem}

This is \cite[Theorem~1.1]{C}. We use its stated conclusion;
we do not interpret its rationalized cohomological construction as
an integral leading-term formula.

\begin{theorem}[Two-primary leading term]\label{classical:two-primary}
For every elliptic curve $A/\Q$ with $s_2(A)\leq1$, the quotient
$Q_A$ of \eqref{eq:normalized-quotient}, with the total real period
and the full-lattice regulator of \eqref{eq:height-convention}, is a
positive rational number, and
\[
 v_2(Q_A)=v_2(\#\Sha(A/\Q)).
\]
\end{theorem}

This is \cite[Theorem~1.1]{B}. Its role here is exactly
the assertion at two, together with positivity and rationality.
After Theorem~\ref{classical:low-corank},
\eqref{eq:kummer-exact} gives $s_p(E)=r(E)$ at every prime. Thus
Theorem~\ref{classical:two-primary} applies to the curve in
Theorem~\ref{thm:main}.

We also use the forward analytic theorem, in order to apply these
conventions to auxiliary twists chosen by analytic nonvanishing.

\begin{theorem}[Analytic rank zero and one]\label{classical:analytic-low-rank}
If an elliptic curve $A/\Q$ has $a(A)\leq1$, then
$r(A)=a(A)$ and the entire group $\Sha(A/\Q)$ is finite.
\end{theorem}

This is the Gross--Zagier--Kolyvagin theorem, using modularity
\cite{GrossZagier,Kolyvagin,BCDT}. By \eqref{eq:kummer-exact} and
Theorem~\ref{classical:two-primary}, $Q_A\in\Q_{>0}$ for every such
$A$. Hence $X_p(A)$ is defined for all the twists used below.
From now on, $p>2$ is fixed, $v=v_p$ on constants, and we abbreviate
$X_p(A)$ to $X(A)$.

\begin{proposition}[Isogeny invariance]\label{classical:isogeny}
Suppose $A$ and $A'$ are isogenous elliptic curves over $\Q$ of
analytic rank at most one. Then
\[
 \frac{Q_A}{\#\Sha(A/\Q)}
       =\frac{Q_{A'}}{\#\Sha(A'/\Q)}.
\]
In particular $X(A)=X(A')$. Over a number field, the corresponding
analytic-to-arithmetic BSD ratio is likewise invariant under isogeny
when the Tate--Shafarevich groups are finite. It is compatible with
products and restriction of scalars.
\end{proposition}

\begin{proof}
These are the arithmetic-volume identities of Cassels and the
restriction-of-scalars comparison of Milne \cite{Cassels,MilneBSD}.
The ratio is formed with the period lattice, full point lattice,
torsion, component factors, and finite Tate--Shafarevich group;
the kernel and cokernel factors of an isogeny cancel in this combined
ratio. The analytic $L$-function is isogeny invariant. Finiteness in
the rational low-rank applications is supplied by
Theorem~\ref{classical:analytic-low-rank}; it passes to the products
and restrictions used here. Thus no conjectural equality of either
ratio with one is required.
\end{proof}

\subsection{Auxiliary quadratic fields}

The pair comparison uses simultaneous nonvanishing for three twists.
We record precisely the density input used for that choice.

\begin{theorem}[Analytic twist densities]\label{classical:twist-density}
For every elliptic curve $A/\Q$, among nonzero squarefree integers
$d$ ordered by $|d|$, the twists $A^d$ of analytic rank zero and one
each have density $1/2$. Consequently the set of $d$ with
$a(A^d)\geq2$ has density zero.
\end{theorem}

We use the theorem at the scope stated in \cite[Theorem~1.2]{Goldfeld}.
Its use below requires no uniform estimate as $A$ varies: one first
chooses one discriminant and then applies the theorem to two fixed
curves. This is the auxiliary-field argument of \cite[Section~2]{C}.

\begin{lemma}[Compatible auxiliary fields]\label{classical:auxiliary-fields}
Let $E/\Q$ have conductor $N$ and analytic rank $r\in\{0,1\}$.
There exist coprime negative odd fundamental discriminants $D,D'<-4$
such that, for
\[
 K=\Q(\sqrt D),\qquad F=\Q(\sqrt{DD'}),\qquad
 L=\Q(\sqrt D,\sqrt{D'}),
\]
the following statements hold.
\begin{enumerate}[label=\textup{(\roman*)}]
\item Every prime dividing $2Np$ splits in both imaginary quadratic
fields, and $\left(\frac{D}{|D'|}\right)=1$.
\item The analytic ranks of $E^D,E^{D'},E^{DD'}$ are respectively
$1-r,1-r,r$.
\item The extension $L/F$ is CM and unramified at every finite place,
and $L$ has only the roots of unity $\{\pm1\}$.
\item One may require $L\cap\Q(E[p])=\Q$ and, when $E$ has CM,
that $L$ not contain its CM field.
\end{enumerate}
Finitely many additional splitting conditions and exclusions of
quadratic fields can be imposed. Both $L(E/K,s)$ and
$L(E^{D'}/K,s)$ have a simple zero at $s=1$.
\end{lemma}

\begin{proof}
Negative odd fundamental discriminants in an admissible fixed
congruence class have positive density. Prescribe splitting at
$2Np$, impose any additional splitting conditions, and exclude a
finite set of quadratic fields. The zero-density exceptional set
in Theorem~\ref{classical:twist-density} cannot exhaust these classes.
We can therefore choose $D$ so that $a(E^D)\leq1$.

Keep $D$ fixed. Choose $D'$ prime to $D$, subject to the same
splitting conditions and, in addition,
$D'\equiv-1\pmod{|D|}$. Together with a sufficiently divisible
modulus supported at $2Np$ and any other prescribed primes, these
conditions contain admissible classes of positive density. Since
$|D'|\equiv1\pmod{|D|}$, quadratic reciprocity gives the required
Jacobi symbol. Apply Theorem~\ref{classical:twist-density} to the
two fixed curves $E$ and $E^D$. The union of their exceptional sets
still has density zero. Thus one can arrange both
$a(E^{D'})\leq1$ and $a(E^{DD'})\leq1$.

For a fundamental discriminant $b$ prime to $N$, the root-number
multiplier is $\chi_b(-N)$. Our splitting and sign prescriptions
make this multiplier $-1$ for $D$ and $D'$, and $+1$ for $DD'$.
The functional equation and the ranks at most one therefore give
the ranks in (ii). The two quadratic factorizations of the
$L$-functions over $K$ prove the last assertion.

The three quadratic subfields of $L$ have discriminants $D,D',DD'$.
The biquadratic discriminant formula yields
\[
 |\operatorname{disc}L|=D^2(D')^2
                       =\operatorname{disc}(F)^2.
\]
The relative discriminant of $L/F$ has norm one, so it is trivial.
The field $F$ is real and $L/F$ is CM. Excluding the finitely many
quadratic subfields responsible for additional roots of unity gives
the last assertion of (iii).

The finite Galois extension $\Q(E[p])/\Q$ has only finitely many
quadratic subfields. A nontrivial intersection with the biquadratic
field $L$ would contain one of the three quadratic subfields of
$L$. Exclude these possibilities when choosing $D$ and $D'$:
after $D$ is fixed, prescribing either $\Q(\sqrt{D'})$ or
$\Q(\sqrt{DD'})$ excludes only finitely many choices of $D'$.
The same argument excludes the CM field. Finite exclusions do not
alter the preceding positive-density argument.
\end{proof}

For the rest of the pair construction, fix such $D,D',K,F,L$ and
write $E'=E^{D'}$ for the companion. In particular $p$ splits
completely in $L$. The rational compact Selmer group over $K$ for
each of $E,E'$ has dimension one. Indeed restriction of scalars is
isogenous to the product of the corresponding two rational twists,
whose analytic ranks sum to one and whose Tate--Shafarevich groups
are finite. A nonzero point in either rank-one line has nonzero
local logarithm at both places above $p$: the kernel of the local
elliptic logarithm on local points is torsion.

\subsection{Heegner points and the height formula}

Put $d_K=|D|$. Let $f\in S_2(\Gamma_0(N))$ be the normalized
rational eigenform attached to $E$. Choose an oriented ideal of
$\mathcal O_K$ of norm $N$ with cyclic quotient, and let $P_X$
be the sum, once over the ideal class group of $K$, of the associated
maximal-order Heegner divisor minus the cusp $\infty$ on $X_0(N)$.
Thus $P_X\in J_0(N)(K)$. There is no division by the class number.

Let $\phi:X_0(N)\to E$ be a rational modular parametrization taking
$\infty$ to zero. We write $\phi_*$ for its map on degree-zero
divisors, and define
\begin{equation}\label{eq:manin-constant}
 P_K=\phi_*P_X,\qquad
 \phi^*\omega_E=c_E f(q)\,\frac{dq}{q},\qquad
 q=e^{2\pi i\tau}.
\end{equation}
The nonzero rational constant $c_E$ is retained exactly; no claim
that it is a $p$-adic unit is used. Write $(f,f)_N$ for the
unnormalized Petersson integral on the ordinary
$\Gamma_0(N)$ quotient, and put
\[
 A_E=\int_{E(\C)}|\omega_E\wedge\overline{\omega_E}|.
\]

Over a number field, $H_{\rm abs}$ denotes the height of
\eqref{eq:height-convention} formed from the absolute logarithmic
$x$-height. It restricts to $H$ over $\Q$. On the Jacobian, use
the Poincare pairing and the orthogonal projector onto the rational
newform factor with its restricted polarization; denote the
resulting absolute diagonal by $H_{\rm abs,f}$.

\begin{theorem}[Split Gross--Zagier normalization]
\label{classical:gross-zagier}
With the preceding conventions, $D$ odd fundamental, $(D,N)=1$,
all level primes split, and $\mathcal O_K^\times=\{\pm1\}$, one has
\begin{align}
 \frac{\sqrt{d_K}}{8\pi^2(f,f)_N}L'(f/K,1)
       &=2H_{\rm abs,f}(P_X),\label{eq:classical-gz}\\
 L'(E/K,1)&=\frac{A_E}{\sqrt{d_K}\,c_E^2}\cdot
                        2H_{\rm abs}(P_K).
                                      \label{eq:classical-gz-elliptic}
\end{align}
In the simple-zero case $P_K$ is nontorsion.
\end{theorem}

The first formula is the conductor-one split Gross--Zagier formula
in the normalization used in \cite[Chapters~I and~V]{GrossZagier};
the conversion is also recorded in \cite[Section~2]{B}. For clarity,
the conversion to the second formula uses the projection formula
for $\phi_*$ and
\[
 (\deg\phi)A_E=8\pi^2c_E^2(f,f)_N.
\]
The relative height over $K$ is $[K:\Q]H_{\rm abs}=2H_{\rm abs}$.
This explains the factor two in both displays. The canonical height
of the origin divisor is $H_{\rm abs}/2$; its Poincare-pairing
diagonal is $H_{\rm abs}$. Thus this convention agrees with
\eqref{eq:height-convention}, without an additional halving.

\subsection{The real arithmetic comparison}

The following equality is needed at odd primes. It is an equality
of real arithmetic ratios before any valuation is taken.

\begin{proposition}[Full-index discrepancy identity]
\label{classical:discrepancy-identity}
Assume $L(E/K,s)$ has a simple zero under the hypotheses of
Theorem~\ref{classical:gross-zagier}. Then
\begin{equation}\label{eq:real-discrepancy}
 \frac{Q_E}{\#\Sha(E/\Q)}
 \frac{Q_{E^D}}{\#\Sha(E^D/\Q)}
 =\frac{[E(K):\Z P_K]^2}
 {c_E^2\prod_\ell c_\ell(E)^2\,\#\Sha(E/K)}.
\end{equation}
The index on the right is the full group index, including torsion.
Consequently, at any prime $p$, if
\[
 n_P=v_p\bigl([E(K)/E(K)_{\rm tors}:\Z\overline P_K]\bigr),
 \qquad \tau_g=v_p(\#E(K)_{\rm tors}),
\]
then
\begin{equation}\label{eq:classical-discrepancy}
 X_p(E)+X_p(E^D)
 =2(n_P+\tau_g)-2v_p(c_E)
             -2\sum_\ell v_p(c_\ell(E))-v_p(\#\Sha(E/K)).
\end{equation}
\end{proposition}

\begin{proof}
Consider $A=\operatorname{Res}_{K/\Q}E_K$. It is isogenous over
$\Q$ to $E\times E^D$. Its point group and Tate--Shafarevich group
are $E(K)$ and $\Sha(E/K)$, and its $L$-function is $L(E/K,s)$.
Its Neron period is $A_E/\sqrt{d_K}$. To see the discriminant
factor, use an integral quadratic basis for $\mathcal O_K$ in the
restriction of the invariant differential; the determinant of the
dual Lie lattice has real Jacobian $1/\sqrt{d_K}$.
The differential stays minimal over $K$: primes of bad reduction
split, and good reduction stays good.

The finite component product for $A$ is
$\prod_\ell c_\ell(E)^2$. At a bad prime, the two split places
each contribute $c_\ell(E)$. At a ramified prime of good reduction,
the special restriction fiber is an extension of the good special
fiber by a vector group and is connected, so there is no extra
component factor. At the remaining good primes the component
factor is one. These period and component calculations are the
restriction-of-scalars arithmetic comparison
\cite{MilneBSD,B}.

Let $P_0$ generate the full free group $E(K)/E(K)_{\rm tors}$,
and write $P_K=mP_0$ modulo torsion. The rank-one regulator for
$A$ is $2H_{\rm abs}(P_0)$; hence
\[
 \frac{2H_{\rm abs}(P_K)}{\Reg_A}=m^2,
 \qquad [E(K):\Z P_K]=|m|\,\#E(K)_{\rm tors}.
\]
Dividing \eqref{eq:classical-gz-elliptic} by the BSD arithmetic
volume just computed gives the right side of
\eqref{eq:real-discrepancy}. Proposition~\ref{classical:isogeny}
identifies the same ratio with the product on its left. Taking
$v_p$ gives \eqref{eq:classical-discrepancy}.
\end{proof}

In \cite[Section~2]{B}, this real identity is subsequently evaluated
at two. The proof above explains why the identity itself applies
at any prime: neither the restriction-of-scalars calculation nor
the Gross--Zagier height identity used a two-primary BSD equality.
What remains in the present paper is to compute the odd-primary
index and Tate--Shafarevich terms on the right of
\eqref{eq:classical-discrepancy} by an exact integral comparison.
